\title{Controlling Text-to-Image Diffusion by Orthogonal Finetuning}

\begin{abstract}
Large-scale text-to-image diffusion models have demonstrated remarkable success in producing photorealistic imagery conditioned on textual input. Nevertheless, steering these models effectively for diverse downstream applications remains a significant, unresolved challenge. In this work, we present a systematic approach for model adaptation—Orthogonal Finetuning (OFT)—designed to refine text-to-image diffusion models for a range of downstream scenarios. In contrast to previous approaches, OFT can rigorously maintain hyperspherical energy, an attribute reflecting pairwise relationships among neurons confined to the unit hypersphere. Our findings indicate that this preservation is pivotal for retaining the semantic generative capacity inherent in text-to-image diffusion models. To further enhance the stability during finetuning, we introduce Constrained Orthogonal Finetuning (COFT), which incorporates a supplementary radius constraint on the hypersphere. In particular, our study focuses on two salient text-to-image finetuning tasks: (1) subject-driven generation, which aims to synthesize subject-specific images from a limited number of reference images paired with a textual prompt; and (2) controllable generation, which extends the model to process and adhere to additional control inputs. Through empirical evaluation, we demonstrate that OFT surpasses existing techniques in both generation fidelity and training efficiency.
\end{abstract}

\section{Introduction}

State-of-the-art text-to-image diffusion models~\cite{saharia2022photorealistic,ramesh2022hierarchical,rombach2022high} have established themselves as highly effective tools for producing high-quality images directed by textual cues. However, relying solely on textual input can lead to ambiguous or imprecise guidance, limiting control over the finer aspects of image synthesis. To mitigate these limitations, our work addresses two principal classes of text-to-image generation problems:

\begin{itemize}[leftmargin=*,nosep]

    \item \textbf{Subject-driven generation}~\cite{ruiz2023dreambooth}: In this scenario, a handful of example images of a particular subject are provided, and the objective is to generate new images featuring the same subject situated within varying contexts as specified by a textual description.
    \item \textbf{Controllable generation}~\cite{zhang2023adding,mou2023t2i}: This task involves the model interpreting an auxiliary control input (such as canny edges or segmentation maps) in addition to a textual prompt, thereby producing images that conform to both sources of guidance.
\end{itemize}

The core challenge in both tasks is to devise methods for finetuning text-to-image diffusion models in a manner that retains the strengths of their pretrained generative capabilities. Essential criteria for a robust finetuning approach include: (1) \emph{training efficiency}, characterized by minimizing both the number of parameters that require updating and the total training time; and (2) \emph{generalizability preservation}, which concerns maintaining the high quality and diversity of outputs generated by the model. Common finetuning practices involve either making small adjustments to existing neuron weights through a modest learning rate (\emph{e.g.}, \cite{ruiz2023dreambooth}) or augmenting the architecture by incorporating small, re-parameterized components (\emph{e.g.}, \cite{hulora2022,zhang2023adding}). However, these approaches, while straightforward, fall short in reliably ensuring that the generative abilities learned during pretraining are conserved. Furthermore, there remains a lack of theoretical guidance for designing effective finetuning protocols and selecting optimal hyperparameters, such as the training duration. One major obstacle is the absence of an established metric to objectively evaluate the extent to which pretrained generative capacity is preserved. Current finetuning schemes tend to operate under the presumption that a smaller Euclidean norm between the parameters of the finetuned and pretrained models reflects better retention of prior knowledge. Consequently, such methods often adopt low learning rates. Nonetheless, although this heuristic may sometimes be valid, we contend that the Euclidean distance alone is insufficient for assessing semantic fidelity. Thus, employing a more structurally informed metric to differentiate between finetuned and pretrained models would contribute substantially to both safeguarding the pretrained abilities and enhancing the stability of finetuning.

Drawing upon empirical findings that hyperspherical similarity effectively captures semantic structure~\cite{liu2017deep,liu2018decoupled,chen2020angular}, our approach leverages hyperspherical energy~\cite{liu2018learning} to represent pairwise relationships between neurons. This metric, defined as the aggregated hyperspherical similarities (such as cosine similarity) across all neuron pairs within a layer, reflects the uniform distribution of neurons on the unit hypersphere~\cite{liu2021learning}. We posit that a finetuned model exhibiting minimal alteration in hyperspherical energy relative to its pretrained counterpart is desirable. An intuitive method would be to incorporate a regularization term during finetuning to retain constant hyperspherical energy, but this approach does not strictly ensure minimization of energy differences. Consequently, we exploit the inherent invariance of hyperspherical energy: applying an identical orthogonal transformation to every neuron leaves pairwise hyperspherical similarity unchanged. Building on this property, we introduce Orthogonal Finetuning~(OFT), a methodology for adapting large text-to-image diffusion models to specific tasks while maintaining their hyperspherical energy unchanged. The essence of OFT is the joint learning of an orthogonal transformation applied to all neurons in a layer, preserving their mutual angles; this can be interpreted as redefining the canonical coordinates for neurons within each layer. Recognizing additionally that the closer the Euclidean distance between the finetuned and original models, the better the pretraining features are retained, we extend OFT with Constrained Orthogonal Finetuning~(COFT). COFT constrains the updated model to remain within a hyperspherical region of predetermined radius centered at the original neuron configuration.

The rationale behind utilizing orthogonal transformations in neuron finetuning draws in part from the principles of the 2D Fourier transform, which enables the decomposition of images into magnitude and phase spectra. Notably, the phase spectrum—encapsulating the angular relationship between inputs and basis elements—retains most of the semantic content. For instance, reconstructing an image using its phase spectrum and a random magnitude spectrum often preserves its core semantics. This observation points to the significance of altering neuron orientations when aiming to achieve semantic modification in generated images—a central objective for both subject-driven and controllable image generation. Nonetheless, if neuron directions are modified without constraint, this freedom can severely compromise the generative capabilities established during pretraining. To address this, we advocate for maintaining the angles between neuron pairs, rooted in the supposition that these angular relationships are essential to a neural network’s knowledge representation. Consequently, we propose the use of layer-wise, shared orthogonal transformations across neurons in each layer, thereby ensuring that hyperspherical energy remains invariant.

Additionally, our work is informed by orthogonal over-parameterized training~\cite{liu2021orthogonal}, a methodology in which classification networks are trained from scratch via orthogonal transformations applied to an initially random network. The rationale here is that a randomly initialized neural network is expected to exhibit low hyperspherical energy, and \cite{liu2021orthogonal} aims to retain this property throughout training, as lower hyperspherical energy is correlated with improved classification generalization~\cite{liu2018learning,lin2020regularizing}. Orthogonal transformations have been shown in \cite{liu2021orthogonal} to provide ample capacity for learning highly generalizable classifiers. In contrast, our focus is on finetuning text-to-image diffusion models, with the objective of enhancing both controllability and generative performance on downstream tasks. It is important to clarify two distinctions between our Orthogonal Finetuning Transformation (OFT) and the method of \cite{liu2021orthogonal}. First, whereas \cite{liu2021orthogonal} minimizes hyperspherical energy, OFT seeks to conserve the hyperspherical energy of the pretrained network, thereby safeguarding its inherent semantic structures during finetuning; indeed, reducing this energy in diffusion models may in fact disrupt pretrained semantic relationships. Second, OFT explicitly strives to limit divergence from the pretrained model—a restriction not present in the approach of \cite{liu2021orthogonal}. Ultimately, effective finetuning requires striking a balance between model adaptability and preservation of pretrained knowledge, and we contend that our OFT method offers a compelling solution in this regard. Our principal contributions are as follows:

\begin{itemize}[leftmargin=*,nosep]

    \item We introduce a new finetuning strategy termed Orthogonal Finetuning (OFT) to enhance controllability in text-to-image diffusion models. To boost stability further, we present a constrained variant of our method that restricts angular deviations from the original pretrained weights.
    \item In contrast to prior finetuning techniques, OFT enables finetuning while rigorously maintaining the hyperspherical energy, a property that, based on our empirical investigations, is key to preserving the generative semantics encoded in the pretrained model.
    \item Our study applies OFT to both subject-driven and controllable image generation tasks. Through extensive experimentation, we show that our method delivers substantial gains over previous approaches with respect to output quality, training convergence rate, and robustness during finetuning. Additionally, OFT is more sample-efficient, requiring significantly fewer training images and epochs to achieve strong performance.
    \item For controllable generation tasks, we propose a novel control consistency metric designed to assess model controllability. The metric works by estimating the intended control signal from generated outputs and comparing it to the original control prompt, providing further evidence for the high degree of controllability enabled by OFT.
\end{itemize}

\section{Related Work}

\textbf{Text-to-image diffusion models}. The field of text-to-image generation has experienced rapid advances~\cite{saharia2022photorealistic,ramesh2022hierarchical,rombach2022high,gu2022vector,nichol2022glide}, largely propelled by breakthroughs in diffusion models~\cite{ho2020denoising,song2021score,song2021denoising,dhariwal2021diffusion} and multimodal representation learning~\cite{su2020vlbert,lu2019vilbert,radford2021learning,wang2021simvlm,li2022blip,li2023blip,alayrac2022flamingo,schuhmann2022laion}. Methods such as GLIDE~\cite{nichol2022glide} and Imagen~\cite{saharia2022photorealistic} implement diffusion processes in the pixel domain; GLIDE integrates a jointly trained text encoder, whereas Imagen relies on a fixed, pretrained text encoder. By contrast, approaches like Stable Diffusion~\cite{rombach2022high} and DALL-E2~\cite{ramesh2022hierarchical} operate in a latent representation space: Stable Diffusion leverages VQ-GAN~\cite{esser2021taming} to derive a visual codebook, and DALL-E2 uses CLIP~\cite{radford2021learning} to align text and image embeddings. Beyond diffusion, the text-to-image problem has also been tackled using generative adversarial networks~\cite{reed2016generative,zhang2017stackgan,xu2018attngan,li2019controllable} and autoregressive models~\cite{ramesh2021zero,ding2021cogview,wu2022nuwa,yuscaling}. Notably, OFT is a model-agnostic finetuning framework and is applicable across different text-to-image diffusion architectures.

\textbf{Subject-driven generation}. To avoid altering the designated subject, approaches such as \cite{avrahami2022blended,nichol2022glide} utilize user-provided masks as auxiliary input conditions. Inversion-based techniques~\cite{choi2021ilvr,dhariwal2021diffusion,ramesh2022hierarchical,gal2022image} can modify surrounding context without affecting the central subject. The method by \cite{hertz2022prompt} enables both local and global edits independently of input masks. Nonetheless, such methods struggle to consistently preserve details related to subject identity. Pivotal Tuning~\cite{roich2022pivotal} addresses this by finetuning a generator around an inverted latent representation, supplemented by a regularization term to maintain identity features. Along similar lines, \cite{nitzan2022mystyle} develops a personalized generative model using a set of individual face images, whereas \cite{casanova2021instance} can synthesize diverse instance variants but at the risk of sacrificing key identifying details. DreamBooth~\cite{ruiz2023dreambooth} offers an alternative by introducing a custom token and employing both reconstruction and class-specific prior losses when finetuning diffusion models with limited subject imagery. OFT builds on the DreamBooth paradigm but departs from standard low-rate finetuning, instead utilizing orthogonal transformations during the finetuning process.

\textbf{Controllable generation}. Image-to-image translation is fundamentally an instance of controllable generation. Historically, most approaches leverage conditional generative adversarial networks~\cite{isola2017image,zhu2017toward,wang2018high,park2019semantic,choi2018stargan} for this task. Recently, diffusion models have also gained popularity for image-to-image translation~\cite{saharia2022palette,voynov2022sketch,wang2022pretraining}. Notably, ControlNet~\cite{zhang2023adding} demonstrates that a pretrained diffusion model can be guided effectively through finetuning, incorporating extra control signals and thus achieving remarkable performance in controllable image generation. In a similar vein, T2I-Adapter~\cite{mou2023t2i} also refines pretrained diffusion models, enhancing their controllability for image synthesis. Building upon the setup proposed in~\cite{zhang2023adding,mou2023t2i}, we employ OFT to further finetune pretrained diffusion models. This leads to superior controllability with less reliance on training data and fewer finetuning parameters. Importantly, OFT does not incur additional inference-time computation.

\textbf{Model finetuning}. Recently, the strategy of finetuning large pretrained models for specialized downstream applications has surged in popularity~\cite{devlin2018bert,brown2020language,he2022masked}. Within this framework, adaptation techniques (\emph{e.g.}, \cite{hulora2022,houlsby2019parameter,pfeiffer2020adapterfusion}) have received significant attention in the field of natural language processing. LoRA~\cite{hulora2022}, which is highly relevant to OFT, models the additive weight update as low-rank during finetuning. Distinctly, OFT introduces layer-shared orthogonal transformation to update neuron weights multiplicatively, and this method guarantees preservation of the pairwise angles among neurons in each layer, resulting in improved stability.

\section{Orthogonal Finetuning}

\subsection{Why Does Orthogonal Transformation Make Sense?}

We begin our analysis by articulating the rationale for employing orthogonal transformations when finetuning text-to-image diffusion models. Our discussion is divided into two fundamental considerations: (1) the motivation for tuning neuron angles (\emph{i.e.}, altering direction), and (2) the reasoning behind selecting orthogonal transformations specifically for this purpose.

To address the first question, we are motivated by empirical findings in \cite{liu2018decoupled,chen2020angular}, which indicate that the angular difference between features effectively reflects the semantic disparity. According to SphereNet~\cite{liu2017deep}, constraining all neurons of a neural network to reside on the surface of a unit hypersphere maintains model expressiveness and can even improve generalization, suggesting that neuron direction suffices to encapsulate the crucial information within data. To more clearly illustrate the role of neuron angles, we set up a toy experiment depicted in Figure~\ref{fig:toy_ae}, in which a conventional convolutional autoencoder is trained on images of flowers. During training, the feature map is obtained via the typical inner product (with $z$ representing the output of the convolution kernel $\bm{w}$ applied to input $\bm{x}$ within a moving window). In the evaluation phase, we consider three alternative methods to reconstruct the feature map: (a) repeating the inner product as in training, (b) leveraging only magnitude information, and (c) relying on angular information. The visual outcomes presented in Figure~\ref{fig:toy_ae} reveal that neuron angular data alone nearly perfectly reconstructs the input images, whereas magnitude provides no informative reconstruction. It should be noted that the cosine-based activation (c) is never used during training, which is based solely on the inner product. This demonstrates that the angles—directions—of the neurons encode the dominant portion of semantic content from the input images. Consequently, adjusting neuron directions should provide the most direct means to manipulate the semantic attributes of images.

In response to the second question, the most direct strategy for adjusting neuron directions involves synchronously rotating or reflecting all neurons within a layer, a process inherently aligned with orthogonal transformation. Although more sophisticated angular operations, such as rotating individual neurons by unique angles, might provide extra flexibility, our findings indicate that orthogonal transformations offer a practical balance between adaptability and structural discipline. Further, as established by \cite{liu2021orthogonal}, orthogonal transformations possess sufficient expressive capacity for neural network training. To empirically reinforce this point, we conducted a regularization analysis under an orthogonality constraint, employing a controllable generation setup as in ControlNet~\cite{zhang2023adding}. The corresponding outcomes, shown in Figure~\ref{fig:ortho}, demonstrate that our canonical OFT, subject to orthogonality, is robust and delivers precise control by epoch 20. Conversely, omitting the orthogonality constraint from OFT results in failure to produce realistic images or to enable any form of control. This experiment underlines the essential contribution of the orthogonality constraint within OFT.

\subsection{General Framework}

Traditional finetuning methods generally employ gradient descent with a modest learning rate, which updates the parameters of a model—or just specific layers—incrementally. The choice of a small learning rate imposes a natural regularization, causing model weights to remain close to those of the pretrained model. Thus, typical finetuning can be viewed as an optimization that, albeit implicitly, minimizes $\thickmuskip=2mu \medmuskip=2mu \| \bm{M} - \bm{M}^0\|$, where $\bm{M}$ represents the parameters after finetuning, and $\bm{M}^0$ are the initial, pretrained weights. While intuitively reasonable, this implicit constraint may be overly permissive, especially for large-scale models. To tighten this restriction, LoRA introduces an explicit low-rank condition on the weight adjustment, requiring that $\thickmuskip=2mu \medmuskip=2mu \text{rank}(\bm{M} - \bm{M}^0)=r'$, where $r'$ is intentionally chosen to be small. In contrast, OFT enforces a different restriction by preserving pairwise neuron similarity through the condition $\thickmuskip=2mu \medmuskip=2mu \| \text{HE}(\bm{M})-\text{HE}(\bm{M}^0) \|=0$, where $\text{HE}(\cdot)$ refers to the hyperspherical energy of a weight matrix. For illustration, consider a dense layer $\thickmuskip=2mu \medmuskip=2mu \bm{W}=\{\bm{w}_1,\cdots,\bm{w}_n\}\in\mathbb{R}^{d\times n}$, where each column $\bm{w}_i\in\mathbb{R}^{d}$ corresponds to neuron $i$ (and $\bm{W}^0$ denotes the pretrained weights). The output $\thickmuskip=2mu \medmuskip=2mu \bm{z}\in\mathbb{R}^n$ of this layer is computed by $\thickmuskip=2mu \medmuskip=2mu \bm{z}=\bm{W}^\top\bm{x}$, where the input is $\thickmuskip=2mu \medmuskip=2mu \bm{x}\in\mathbb{R}^d$. Under OFT, the training objective minimizes the difference in hyperspherical energy between the finetuned and pretrained models:

\begin{equation}\label{eq:oft_principle}
\footnotesize
\min_{\bm{W}} \left\| \text{HE}(\bm{W})-\text{HE}(\bm{W}^0)\right\|~~\Leftrightarrow~~\min_{\bm{W}} \bigg{\|} \sum_{i\neq j} \|\hat{\bm{w}}_i-\hat{\bm{w}}_j\|^{-1}-\sum_{i\neq j} \|\hat{\bm{w}}^0_i-\hat{\bm{w}}^0_j\|^{-1}\bigg{\|}
\end{equation}

Here, the normalized neuron is given by $\thickmuskip=2mu \medmuskip=2mu \hat{\bm{w}}_i:=\bm{w}_i/\|\bm{w}_i\|$, and for a dense layer $\bm{W}$, the hyperspherical energy is defined as $\thickmuskip=2mu \medmuskip=2mu \text{HE}(\bm{W}):= \sum_{i\neq j} \|\hat{\bm{w}}_i-\hat{\bm{w}}_j\|^{-1}$. It is apparent that Eq.~\eqref{eq:oft_principle} achieves a minimum of zero. This minimum is reached provided that $\bm{W}$ and $\bm{W}^0$ differ by no more than a rotation or reflection, that is, when $\thickmuskip=2mu \medmuskip=2mu \bm{W}=\bm{R}\bm{W}^0$, with $\thickmuskip=2mu \medmuskip=2mu \bm{R}\in\mathbb{R}^{d\times d}$ an orthogonal transformation (rotation or reflection determined by determinant $1$ or $-1$). The core idea of OFT is thus to restrict finetuning to learning an orthogonal transformation

\begin{equation}\label{eq:oft_general}
    \footnotesize
    \bm{z}=\bm{W}^\top\bm{x}=(\bm{R}\cdot\bm{W}^0)^\top\bm{x},~~~\text{s.t.}~\bm{R}^\top\bm{R}=\bm{R}\bm{R}^\top=\bm{I}
\end{equation}

Here, $\bm{W}$ represents the weight matrix obtained via OFT finetuning, while $\bm{I}$ indicates the identity matrix. An overview of OFT can be found in Figure~\ref{fig:block}. Analogous to LoRA’s use of zero initialization, it is essential that OFT performs finetuning on top of the pretrained model exactly from $\bm{W}^0$. This is achieved by initializing the orthogonal matrix $\bm{R}$ as the identity, ensuring the initial weights match those of the pretrained model. Maintaining the orthogonality constraint on $\bm{R}$ can be accomplished using differentiable orthogonalization methods proposed in prior works~\cite{liu2021orthogonal,lezcano2019cheap}. Later, we elaborate on computationally efficient approaches for enforcing this orthogonality.

\subsection{Efficient Orthogonal Parameterization}\label{sect:eff_ortho}

Typical orthogonalization procedures such as the differentiable Gram-Schmidt algorithm tend to be computationally prohibitive in practical scenarios~\cite{liu2021orthogonal}. To improve efficiency, we utilize the Cayley parameterization for constructing orthogonal matrices. Specifically, the orthogonal matrix is formulated as $\thickmuskip=2mu \medmuskip=2mu \bm{R}=(\bm{I}+\bm{Q})(I-\bm{Q})^{-1}$, where $\bm{Q}$ is skew-symmetric, i.e., $\thickmuskip=2mu \medmuskip=2mu \bm{Q}=-\bm{Q}^\top$. One trade-off of this method is that the Cayley parameterization only generates orthogonal matrices with determinant $1$, restricting $\bm{R}$ to lie in the special orthogonal group. Empirically, this constraint does not compromise performance. Nonetheless, even with the Cayley transform, using a full-size $\bm{R}$ becomes inefficient as dimensionality $d$ grows. To improve parameter efficiency, we introduce a block-diagonal structure for $\bm{R}$, composed of $r$ independent blocks. This results in the following structure:

\begin{equation}
    \footnotesize
    \bm{R}=\text{diag}(\bm{R}_1,\bm{R}_2,\cdots,\bm{R}_r)=\begin{bmatrix}
    \bm{R}_1\in \text{O}(\frac{d}{r}) &  &\\
    &\ddots & \\
    & & \bm{R}_r\in \text{O}(\frac{d}{r})
    \end{bmatrix}\in \text{O}(d)
\end{equation}

where $\text{O}(d)$ stands for the group of $d$-dimensional orthogonal matrices, with $\thickmuskip=2mu \medmuskip=2mu \bm{R}\in\mathbb{R}^{d\times d}$ and each block $\thickmuskip=2mu \medmuskip=2mu \bm{R}_i\in\mathbb{R}^{d/r\times d/r},\forall i$ is an orthogonal matrix as well. In the special case where $\thickmuskip=2mu \medmuskip=2mu r=1$, the block-diagonal structure reduces to an unconstrained orthogonal matrix. An orthogonal matrix of dimension $\thickmuskip=2mu \medmuskip=2mu d\times d$ contains $\thickmuskip=2mu \medmuskip=2mu d(d-1)/2$ independent parameters, yielding a computational complexity of $\mathcal{O}(d^2)$. In the $r$-block diagonal scenario, the parameter count is given by $\thickmuskip=2mu \medmuskip=2mu d(d/r-1)/2$, so the complexity drops to $\thickmuskip=2mu \medmuskip=2mu \mathcal{O}(d^2/r)$. An additional reduction can be achieved by sharing the same block matrix across all blocks, i.e., enforcing $\thickmuskip=2mu \medmuskip=2mu\bm{R}_i=\bm{R}_j$ for all $i\neq j$. This leads to parameter complexity scaling as $\mathcal{O}(d^2/r^2)$. Importantly, these approaches enhance parameter efficiency while retaining the orthogonality of $\bm{R}$, ensuring that the hyperspherical energy is strictly maintained.

We proceed to compare OFT and LoRA regarding parameter usage. For LoRA with a rank-reduction parameter $r'$, the number of learnable weights is $\thickmuskip=2mu \medmuskip=2mu r'(d+n)$. If we parameterize both $r$ and $r'$ as fractions of the model dimension $d$ (for instance, setting $\thickmuskip=2mu \medmuskip=2mu r=r'=\alpha d$ with $\thickmuskip=2mu \medmuskip=2mu 0<\alpha\leq 1$), LoRA’s complexity is then $\mathcal{O}(d^2+dn)$, whereas OFT achieves a complexity of only $\mathcal{O}(d)$. For illustration, consider a $128\times128$ weight matrix: with $\thickmuskip=2mu \medmuskip=2mu r'=8$, LoRA involves $2,048$ trainable parameters, but OFT requires only $960$ parameters for $\thickmuskip=2mu \medmuskip=2mu r=8$, without applying block sharing.

\subsection{Constrained Orthogonal Finetuning}

To further regulate the adaptability of OFT, it is possible to confine the finetuned parameters within a restricted vicinity around the pretrained weights. Constrained Orthogonal Finetuning (COFT) does so by applying the following operation:

\begin{equation}\label{eq:coft_general}
    \footnotesize
    \bm{z}=\bm{W}^\top\bm{x}=(\bm{R}\cdot\bm{W}^0)^\top\bm{x},~~~\text{s.t.}~\bm{R}^\top\bm{R}=\bm{R}\bm{R}^\top=\bm{I},~\norm{\bm{R}-\bm{I}}\leq \epsilon
\end{equation}

This setup enforces both the orthogonality of $\bm{R}$ and bounds its deviation from the identity matrix by $\epsilon$. The orthogonal constraint can be addressed by leveraging the Cayley parameterization as discussed in Section~\ref{sect:eff_ortho}. Nonetheless, imposing the $\epsilon$-deviation limit in the context of Cayley-parameterized matrices poses non-trivial challenges. To explore this, we utilize the Neumann series to approximate the Cayley transform $\thickmuskip=2mu \medmuskip=2mu \bm{R}=(\bm{I}+\bm{Q})(I-\bm{Q})^{-1}$, arriving at the expansion $\thickmuskip=2mu \med

series is convergent in the operator norm. Consequently, it is possible to reposition the constraint $\thickmuskip=2mu \medmuskip=2mu\|\bm{R}-\bm{I}\|\leq\epsilon$ within the Cayley transform. This leads to an alternative but equivalent requirement: $\thickmuskip=2mu \medmuskip=2mu \|\bm{Q}-\bm{0}\|\leq\epsilon'$, where $\bm{0}$ stands for the zero matrix and $\epsilon'$ is a distinct tolerance hyperparameter from $\epsilon$. The constraint on $\bm{Q}$ can be straightforwardly managed using projected gradient descent methods. To obtain an orthogonal matrix $\bm{R}$ initialized at identity, one can simply set the initial value of $\bm{Q}$ to the zero matrix. The COFT approach can thus be interpreted as enforcing two concrete constraints: maintaining a low hyperspherical energy difference, and limiting divergence from the original pretrained weights. While the latter is often enforced implicitly by standard finetuning approaches, COFT formalizes it as an explicit constraint. Although OFT demonstrates strong performance, our experiments indicate that COFT achieves greater finetuning stability, owing to its explicit regularization on model deviation. Figure~\ref{fig:eps_coft} illustrates the influence of the parameter $\epsilon$ on COFT’s behavior. It is apparent that $\epsilon$ modulates the scope of finetuning adaptation: larger values of $\epsilon$ result in COFT-finetuned models that closely approximate those from OFT, whereas smaller values produce models that retain more properties of the original pretrained text-to-image diffusion model.

\subsection{Re-scaled Orthogonal Finetuning}

We introduce a straightforward modification to standard OFT by learning an individual scaling parameter for every neuron. The rationale stems from the observation that scaling neuron outputs does not impact the hyperspherical energy, since the normalization process negates any magnitude effects. The resulting forward computation takes the form: $\thickmuskip=2mu \medmuskip=2mu\bm{z}=(\bm{R}\bm{W}^0\bm{D})^\top\bm{x}$\footnote{\textbf{Errata}: The NeurIPS camera-ready version inadvertently reported the forward pass as $\thickmuskip=2mu \medmuskip=2mu\bm{z}=(\bm{D}\bm{R}\bm{W}^0)^\top\bm{x}$. However, the original code implementation is accurate, so the findings in Appendix~\ref{app:rescaled} are correct.} with $\thickmuskip=2mu \medmuskip=2mu \bm{D}=\text{diag}(s_1,\cdots,s_n)\in\mathbb{R}^{n\times n}$ representing a diagonal matrix of learnable parameters $s_1,\cdots,s_n$ where all entries are strictly positive. Unlike the original OFT forward equation Eq.~\eqref{eq:oft_general} where $\bm{R}$ alone is trainable, both $\bm{R}$ and the diagonal scaling matrix $\bm{D}$ are now optimized. This re-scaled OFT variant enhances the flexibility of OFT with negligible increase in the total parameter count. For experimental clarity, we focus on basic OFT to highlight the benefit of the orthogonal mapping, though we note empirically that the re-scaled version typically yields even better outcomes (see Appendix~\ref{app:rescaled}).

\section{Intriguing Insights and Discussions}

\textbf{OFT’s applicability across network designs}. OFT is, in principle, adaptable to any neural network architecture. In Transformers, LoRA is predominantly employed on the attention layers~\cite{hulora2022}. To ensure a direct comparison to LoRA, our experimental setup restricts OFT application to attention weights. In addition to dense (fully connected) layers, OFT is well-suited for tuning convolutional layers, with its block-diagonal matrix $\bm{R}$ acquiring a meaningful structure in the convolutional setting—an advantage not shared by LoRA. By configuring the number of blocks to match the input channel count, each block applies an orthogonal transformation uniquely to a single neuron channel, much like learning depthwise convolution filters~\cite{chollet2017xception}. Alternatively, sharing the blocks of $\bm{R}$ causes OFT to impose a unified orthogonal mapping to all channels simultaneously. Appendix~\ref{app:convolution} presents an initial exploration of OFT for finetuning convolutional layers.

\textbf{Connection to LoRA}. LoRA mitigates significant changes in the pretrained weight matrix by incorporating a low-rank matrix. In comparison, OFT imposes an orthogonal (full-rank) transformation on the pretrained weights, thus safeguarding the pretrained information from being disrupted by the transformation. The forward computation of OFT can be reformulated as $\thickmuskip=2mu \medmuskip=2mu \bm{z}=(\bm{R}\bm{W}^0)^\top\bm{x}=(\bm{W}^0+(\bm{R}-\bm{I})\bm{W}^0)^\top\bm{x}$, where the term $\thickmuskip=2mu \medmuskip=2mu (\bm{R}-\bm{I})\bm{W}^0$ plays a role similar to the low-rank update in LoRA. Provided $\bm{W}^0$ possesses full rank, OFT also achieves a low-rank modification whenever $\thickmuskip=2mu \medmuskip=2mu \bm{R}-\bm{I}$ itself is of low rank. LoRA controls its rank via the hyperparameter $r'$, and similarly, OFT uses a block-diagonal size $r$ to constrain its trainable parameters. Fundamentally, LoRA and OFT illustrate two divergent forms of parameter-efficient adaptation: LoRA reduces parameters by leveraging low-rankness, whereas OFT achieves efficiency by utilizing the sparsity pattern, specifically block-diagonal orthogonality.

\textbf{Why OFT converges faster}. Firstly, as Figure~\ref{fig:toy_ae} illustrates, the optimal approach to change the representational semantics is through altering the directions of neurons—precisely the mechanism employed by OFT. Secondly, OFT can be interpreted as adapting neurons while constrained to a smooth hypersphere, resulting in an optimization landscape with more favorable properties. This intuition aligns with empirical observations in \cite{liu2021orthogonal}.

\textbf{Why not minimize hyperspherical energy}. Unlike the approach in \cite{liu2021orthogonal}, our method does not pursue minimizing hyperspherical energy. In classification settings, maintaining non-redundant neurons is advantageous. Minimizing hyperspherical energy enforces even spacing of neurons on the hypersphere, but such a criterion is ill-suited to finetuning tasks since it risks erasing the useful information learned during pretraining.

\textbf{Balancing flexibility and regularity during finetuning}. Our study reveals a fundamental balance that must be struck between the adaptability (flexibility) and the consistency (regularity) of finetuning approaches. Conventional finetuning achieves maximal adaptability but at the cost of limited robustness and a propensity for model collapse. OFT, despite its simplicity, achieves a more optimal compromise by maintaining pairwise neuron angles, thus managing both flexibility and regularization. The block-diagonal approach can be interpreted as imposing even stricter regularization constraints on the orthogonal matrix.

\textbf{Zero inference-time computational cost}. In contrast to methods such as ControlNet, the OFT approach does not result in any inference-time computational burden on the finetuned model. During inference, we can directly compose the learned orthogonal transformation matrix $\bm{R}$ with the original pretrained weights $\bm{W}^0$, producing an updated weight matrix $\thickmuskip=2mu \medmuskip=2mu \bm{W}=\bm{R}\bm{W}^0$. Consequently, the inference latency remains unchanged compared to the base pretrained model.

\section{Experiments and Results}

\textbf{Experimental protocol}. For our empirical analysis, Stable Diffusion v1.5~\cite{rombach2022high} is employed as the foundational text-to-image pretrained model. To ensure impartiality, output samples from each approach are randomly selected. Procedures for subject-driven synthesis generally replicate the pipeline of DreamBooth~\cite{ruiz2023dreambooth}. For tasks involving controlled image generation, we largely adhere to ControlNet~\cite{zhang2023adding} and T2I-Adapter~\cite{mou2023t2i}. In order to provide a direct comparison with LoRA, we restrict OFT or COFT to the same model layers used by LoRA. Additional findings and implementation specifics can be found in Appendix~\ref{app:settings}.

\subsection{Subject-driven Generation}

\textbf{Experiment details}. DreamBooth~\cite{ruiz2023dreambooth} and LoRA~\cite{hulora2022} serve as comparative baselines. All tested methods employ the loss formulation specified in DreamBooth. When applying DreamBooth and LoRA, we closely follow the original papers’ recommended hyperparameter configurations. Further experimental outcomes are available in Appendix~\ref{app:settings},\ref{app:coft_oft},\ref{app:more_qualitative},\ref{app:failure}.

\textbf{Finetuning stability and convergence}. To begin, we assess the stability during finetuning as well as the convergence rate across DreamBooth, LoRA, OFT, and COFT. Refer to Figure~\ref{fig:teaser} and Figure~\ref{fig:db_convergence} for empirical results. It is evident that both COFT and OFT demonstrate robust and stable finetuning for the diffusion model. At 400 training steps, DreamBooth and the OFT-based approaches provide effective controllability, whereas LoRA is unable to maintain subject identity. Upon reaching 2000 iterations, DreamBooth exhibits mode collapse by producing unrealistic outputs, and LoRA cannot render a yellow shirt, generating yellow fur instead. In comparison, OFT and COFT continue to deliver stable and coherent image control. These findings substantiate the ability of our OFT framework to combine rapid convergence with stability for subject-driven image generation. A notable advantage is the reduced sensitivity to the total number of finetuning iterations, simplifying the process and mitigating the need for per-subject iteration adjustment. Both OFT and COFT support using larger iteration counts by default. For COFT, with an adequately chosen $\epsilon$, both learning rate and number of iterations require minimal tuning.

\textbf{Quantitative comparison}. Adopting the evaluation strategy from \cite{ruiz2023dreambooth}, we present a quantitative analysis on subject fidelity (using DINO~\cite{caron2021emerging} and CLIP-I~\cite{radford2021learning}), adherence to text prompts (CLIP-T~\cite{radford2021learning}), and output diversity (LPIPS~\cite{zhang2018unreasonable}). For CLIP-I, we calculate the mean pairwise cosine similarity between CLIP feature vectors for real and generated images; DINO operates similarly to CLIP-I, but utilizes ViT S/16 DINO embeddings. CLIP-T is determined by averaging the cosine similarity between textual prompts and their corresponding generated images using CLIP. The LPIPS metric measures average cosine similarity in feature space for images of the same subject generated from identical prompts. According to Table~\ref{table:dreambooth_final}, both COFT and OFT significantly exceed DreamBooth and LoRA in DINO and CLIP-I subject fidelity, while also delivering equal or improved results in text prompt fidelity and diversity. To address randomness, each algorithm is independently finetuned on the same pretrained model 30 times with distinct random seeds. Collectively, these results illustrate that OFT is not only more stable and convergent but also achieves more reliable final metrics across experiments.

\textbf{Qualitative comparison}. To provide a clearer picture of the advantages offered by OFT, we present a selection of randomly sampled cases for subject-driven generation in Figure~\ref{fig:dreambooth_final}. In order to ensure fairness across methods, all samples are created from the same finetuned base model under each approach, with text prompts kept fixed so that no prompt tuning is done to favor individual outputs. The model with the optimal validation CLIP score is chosen for each method. Examining Figure~\ref{fig:dreambooth_final}, we notice that both OFT and COFT display strong abilities to maintain the semantic content of the subject, whereas LoRA frequently struggles with subject identity consistency (for instance, LoRA entirely loses the subject in the bowl scenario). Additionally, OFT and COFT exhibit superior responsiveness to text guidance, in contrast to DreamBooth, which, despite preserving subject identity, often cannot faithfully generate content as specified by the text prompt (e.g., see the initial row for the bowl case). These observations indicate that OFT achieves a better balance between controllability via prompts and fidelity to subject identity. Notably, OFT is robust to the number of finetuning iterations and yields high-quality results even with extensive training; this level of stability is not seen with DreamBooth or LoRA. Further examples are included in Appendix~\ref{app:more_qualitative}. Human studies, described in Appendix~\ref{app:human}, also support the superior performance of OFT.

\subsection{Controllable Generation}

\textbf{Settings}. Our experimental setup includes ControlNet~\cite{zhang2023adding}, T2I-Adapter~\cite{mou2023t2i}, and LoRA~\cite{hulora2022} as comparison baselines. We focus on three demanding controllable generation benchmarks: translating Canny edge maps to images~(C2I) on the COCO dataset~\cite{lin2014microsoft}, converting segmentation maps to images~(S2I) on ADE20K~\cite{zhou2017scene}, and producing faces from landmarks~(L2F) on CelebA-HQ~\cite{karras2018progressive,xia2021tedigan}. Each technique is used to finetune Stable Diffusion~(SD) v1.5 for 20 epochs on all three datasets. Supplementary qualitative samples can be found in Appendix~\ref{app:more_qualitative},\ref{app:more_control_task},\ref{app:failure}.

\textbf{Convergence}. We assess the convergence properties of ControlNet, T2I-Adapter, LoRA, and COFT within the L2F task through both quantitative and qualitative analyses. For quantitative measurement, we utilize the mean $\ell_2$ distance between the intended control face landmarks and the generated face landmarks. Figure~\ref{fig:face_convergence} presents the progression of the face landmark error as each model is finetuned across various epochs. Notably, both COFT and OFT demonstrate much more rapid convergence relative to the other approaches; for example, LoRA requires 20 epochs to reach the same error level that OFT attains after just 8 epochs. Importantly, OFT and COFT maintain a comparable parameter count to LoRA (and substantially fewer than ControlNet), yet they achieve considerably faster convergence than the alternative models. Furthermore, Figure~\ref{fig:teaser} corroborates the data efficiency of OFT; the rightmost illustration reveals that OFT is capable of effective convergence using merely 5\% of the ADE20K dataset, outperforming ControlNet and LoRA in terms of required data. For qualitative comparison, we limit our attention to OFT, COFT, and ControlNet, given that ControlNet attains landmark errors nearest to those of our approach. The qualitative results illustrated in Figure~\ref{fig:face_convergence_qual} indicate that both OFT and COFT converge consistently, with the predicted face pose increasingly aligning to the reference control landmarks over time. In comparison, ControlNet exhibits a more abrupt increase in controllability after the eighth epoch—an abrupt transition also reported in \cite{zhang2023adding}. This enhanced convergence stability inherent to the OFT framework translates to improved ease of practical application, as OFT's training trajectory is more gradual and predictable, facilitating hyperparameter selection.

\textbf{Quantitative comparison}. To assess the effectiveness of controllable generation, we propose a metric designed to quantify control consistency. This approach involves extracting the control signal from each synthesized image and directly comparing it against the corresponding original input signal. For the C2I setting, we report both the IoU and F1 scores, whereas the S2I scenario relies on the mean IoU along with mean and overall accuracy. For the L2F setting, we measure the mean $\ell_2$ distance between the predicted and actual control landmarks. Comprehensive explanations about these consistency metrics can be found in Appendix~\ref{app:settings}. All baseline methods were tested using their most optimal hyperparameter configurations. Table~\ref{table:control_gen} summarizes the outcomes, indicating that both OFT and COFT offer substantially improved and more reliable control than competing approaches. We note that adapter-based strategies (such as T2I-Adapter and ControlNet) tend to have slower convergence and poorer ultimate performance. LoRA demonstrates better results on S2I compared to ControlNet, but underperforms on both C2I and L2F. Generally, despite careful hyperparameter tuning, LoRA appears to have a lower performance upper bound. In contrast, our OFT framework does not appear to reach this plateau; performance continues to improve as the count of trainable parameters increases. Notably, our work introduces a rigorous quantitative protocol for evaluating controllable generation, enabling accurate assessment of the fine-tuned models’ control capabilities—subject to the fidelity of pre-trained segmentation and detection backbones.

\textbf{Qualitative comparison}. We further assess OFT and COFT through qualitative analysis against prevailing state-of-the-art approaches, specifically ControlNet, T2I-Adapter, and LoRA. As depicted by the random sample images in Figure~\ref{fig:control_final}, OFT and COFT consistently produce images that are not only realistic and visually convincing but also exhibit precise controllability. Focusing on the S2I task, it is evident that LoRA fails to adhere to the provided segmentation maps during image synthesis, while ControlNet, OFT, and COFT successfully regulate the output to correspond with input constraints. Notably, in comparison to ControlNet, OFT and COFT are capable of generating images with higher realism, enhanced detail, and more coherent geometry, all while utilizing a significantly reduced number of parameters. For the C2I task, OFT and COFT effectively generate semantically relevant images from coarse Canny edge inputs, whereas T2I-Adapter and LoRA show notably inferior performance. During the L2F task, our approach achieves the highest accuracy in pose control for generated facial images, even when faced with complex facial orientations. Collectively, across all three control scenarios, OFT and COFT surpass state-of-the-art benchmarks in qualitative image outcomes, substantiating the superior controllability of the OFT framework. Additional qualitative results for the three tasks can be found in Appendix~\ref{app:control_exp}, and we further demonstrate OFT’s capability across a broader range of control applications—including dense pose to human body, sketch to image, and depth to image—in Appendix~\ref{app:more_control_task}.

\section{Concluding Remarks and Open Problems}

Recognizing that the angular relationships among neurons are pivotal to encoding visual semantics, we introduce a straightforward yet robust finetuning strategy—orthogonal finetuning—for the purpose of refining text-to-image diffusion models. Our approach specifically addresses two key applications: subject-driven generation and controllable generation. When compared to prevailing techniques, OFT offers superior controllability and stability during finetuning, all while requiring a smaller parameter set. Furthermore, OFT incurs no additional computational burden during inference, making it suitable for efficient real-world deployment.

OFT also opens up several avenues for future investigation. Firstly, orthogonality in OFT is ensured through Cayley parametrization, which necessitates computing a matrix inverse—a step that slightly impedes the scalability of the approach. While we partially mitigate this issue by employing block diagonal parametrization, finding a differentiable method to accelerate this matrix inversion process remains unresolved. Secondly, OFT exhibits notable promise in compositionality: the orthogonal matrices arising from different OFT finetuning procedures can be multiplied, yielding another orthogonal matrix. It is an open question whether the resulting set of orthogonal matrices effectively retains the knowledge acquired from each downstream task. Lastly, OFT’s efficiency in parameter utilization depends heavily on the use of block diagonal structures, which inherently impose extra bias and constrain adaptability. Identifying strategies that further enhance parameter efficiency with minimal bias and greater flexibility is an important direction for future work.

\end{document}